\section{Related Work}
\label{2}
\subsection{Machine Unlearning}
Machine unlearning has evolved through several paradigms, each with efficiency trade-offs. Zhao \textit{et al.}~\cite{zhao2024continual} proposed GS-LoRA using group sparse regularization with LoRA, but bounded gradient ascent slows convergence. Kurmanji \textit{et al.}~\cite{kurmanji2023towards} introduced SCRUB with bidirectional logit manipulation, yet unbounded KL maximization creates convergence difficulties. Recognizing uniform parameter treatment as suboptimal, Fan \textit{et al.}~\cite{fan2023salun} proposed SalUn, which uses gradient-based weight saliency to separate forget and retain parameters, but entangled updates during training violate this mutual exclusivity and random labels can be reverse-mapped. Chen \textit{et al.}~\cite{chen2023boundary} proposed Boundary Unlearning, adding a shadow node and pruning the head for irreversible forgetting, but all class spaces must migrate toward the shadow, causing slow convergence. Sun \textit{et al.}~\cite{sun2024forget} introduced Forget Vectors, tuning noise to suppress forget classes without retraining. However, it remains vulnerable to reverse-engineering attacks that violate GDPR.

Recent methods adopt optimization techniques from related fields. Shi \textit{et al.}~\cite{shi2025mcu} proposed MCU, using Bézier curves between initial and clean models, but incurs dual-optimization overhead. Huang \textit{et al.}~\cite{huang2024learning} introduced LTU, a three-phase meta-learning framework that is computationally expensive. Zhou \textit{et al.}~\cite{zhou2025decoupled} and Spartalis \textit{et al.}~\cite{spartalis2025lotus} proposed DELETE and LoTUS, both requiring frozen teachers that double GPU memory. Wu \textit{et al.}~\cite{wu2025munba} proposed MUNBa, a Nash bargaining framework whose dual-gradient updates add per-step optimization overhead.

However, few existing methods explicitly consider biometric unlearning on edge devices or are validated across sequential forget cycles, and GS-LoRA~\cite{zhao2024continual} targets continual forgetting only for classification.

\subsection{Biometric Identifiers}
\textbf{Iris recognition:}
Kadhim \textit{et al.}~\cite{kadhim2025deep} proposed \textit{SSLDMNet-Iris}, combining parallel multi-scale convolutional branches with GRU branches for synchronized temporal modeling and Fisher's Linear Discriminant for feature separability. Sun \textit{et al.}~\cite{sun2024irisformer} proposed \textit{IrisFormer}, a Vision Transformer with relative positional encoding for rotation invariance, pixel-shift augmentation and random token masking for occlusion robustness, and patch-wise sequential matching. Luo \textit{et al.}~\cite{luo2025novel} proposed a ResNet-50 iris network with a Lightweight Attention Module (LAM) integrating 1D convolution channel attention and cascaded $3\times 3$ spatial attention, trained with an enhanced triplet loss that adaptively mines hard samples.

\textbf{Fingerprint recognition:}

Yu \textit{et al.}~\cite{yu2024partial} proposed \textit{PFSimNet}, a deep learning partial fingerprint matching network that computes a 3D similarity matrix between local feature descriptors from two branches to simulate traditional brute-force descriptor matching, combined with a two-step pre-training flow using an alignment network that forces impostor pairs to share similar orientation fields so the model learns subtle identity-discriminative features. Guan \textit{et al.}~\cite{guan2024joint} proposed \textit{JIPNet}, a CNN-Transformer hybrid network that jointly estimates identity verification and relative pose for partial fingerprint pairs by exploiting their coupling relationship, using shared-weight convolutional blocks for independent feature extraction followed by interleaved self  and cross attention transformers for interrelated feature interaction, along with a lightweight fingerprint enhancement pre-training task to improve feature robustness across modalities.

\textbf{Face recognition:}
Zhong \textit{et al.}~\cite{zhong2021face} proposed \textit{Face Transformer}, a ViT-based face recognition model that uses sliding overlapping patches instead of non-overlapping ones to preserve inter-patch information, achieving performance comparable to CNN backbones with similar parameters and MACs when trained on MS-Celeb-1M. Wu \textit{et al.}~\cite{wu2025privacy} proposed \textit{MEM-FR}, a privacy-preserving face recognition system that performs optical convolution via a mask-encoded microlens array before signals reach the sensor, paired with a backend FaceNet and a dual cosine similarity loss, achieving 95.0\% simulation and 92.3\% physical accuracy on LFW while improving spatial resolution, light throughput, and field of view over the prior LOEN system.
